\section{Mato Grosso: Contexto e Dados}
\label{sec:dados}

\subsection{Unidade territorial, período e biomas}

O painel proposto terá como unidade a combinação entre território municipal comparável e ano. O período candidato de 2001--2024 preserva a escolha registrada na preparação da qualificação, mas permanece condicionado à disponibilidade conjunta das variáveis. O inventário das bases, realizado em setembro de 2026, ainda não constitui uma amostra analítica: há 142 códigos na malha municipal utilizada e 141 em fontes como MapBiomas e Pesquisa da Pecuária Municipal. A criação de Boa Esperança do Norte e seus desmembramentos exigem compatibilização com Sorriso e Nova Ubiratã antes de qualquer junção.

A preferência será por unidades territoriais estáveis ao longo da amostra, mediante agregação de municípios afetados por desmembramentos e documentação da correspondência histórica. Não se atribuirão retrospectivamente dados produtivos à malha atual por simples proporcionalidade de área. Para rasters de cobertura, a agregação poderá ser refeita sobre a geografia harmonizada; para as estatísticas municipais, deverá ser respeitada a unidade de divulgação de cada período. Cada exclusão ou agregação será registrada, incluindo a perda de observações por ausência de informação.

A classificação ambiental utilizará a delimitação de biomas do IBGE de 2019 \cite{IBGEBiomas2019}. Como regra candidata, municípios com pelo menos 80\% de sua área em um bioma serão classificados por predominância; os demais constituirão um grupo misto. Serão examinados limiares de 70\% e 90\% e, alternativamente, as frações contínuas de área em cada bioma. No inventário atual, antes da harmonização histórica, o limiar de 80\% identifica apenas dois municípios predominantemente pantaneiros. Esse número recomenda cautela: resultados separados para esse grupo dependerão de suporte suficiente e poderão permanecer descritivos.

Bioma e regime jurídico serão tratados separadamente. A Amazônia Legal inclui Mato Grosso e não se confunde com o bioma Amazônia. A Lei n.\ 12.651/2012 já estabelece diferenciações na Reserva Legal por localização e formação vegetal, além de condições e exceções específicas \cite[arts.~3 e 12]{Brasil2012}. Portanto, a pesquisa não partirá da existência de uma regra ambiental absolutamente uniforme entre biomas. A questão diz respeito à interação entre regras, implementação e condições econômicas locais. A redação consolidada atual não será retroagida a 2001; a análise institucional exigirá a reconstrução da vigência das normas em cada período.

\subsection{Vegetação e conversões brutas}

A fonte ambiental principal será a coleção 10.1 do MapBiomas, já obtida para 1985--2024. O projeto oferece estatísticas por classe e mapas anuais de cobertura e uso da terra \cite{MapBiomas2026}. A existência de coleções posteriores não implica substituição automática da versão local, cuja identificação deve acompanhar cada produto empírico. As transições serão calculadas entre pixels observados de uma mesma coleção, com legenda, resolução e máscaras consistentes.

Sejam $N_{i,t-1}$ a área de vegetação nativa observável no início do intervalo e $D_{it}^{j}$ a área que passa de vegetação nativa em $t-1$ ao destino $j$ em $t$. A variável dependente será:
\begin{equation}
 C_{it}^{j}=100\,\frac{D_{it}^{j}}{N_{i,t-1}},
 \qquad j\in\{\text{agricultura},\text{pastagem},\text{agropecuária}\}.
 \label{eq:conversao}
\end{equation}
A taxa será expressa em porcentagem da vegetação disponível; diferenças entre taxas serão expressas em pontos percentuais. Agricultura e pastagem terão destinos mutuamente exclusivos; a conversão agropecuária total poderá incluir classes de uso agropecuário não desagregadas, que serão apresentadas separadamente. O denominador deve corresponder aos pixels nativos com observação válida nos dois anos. Territórios sem vegetação inicial não terão taxa definida e serão descritos à parte.

A taxonomia candidata do estoque local reúne formações florestais, savânicas, campestres e outras classes naturais compatíveis com a legenda. A inclusão de cada classe será validada antes da estimação. No Pantanal, em particular, a presença de uso pecuário em campo natural não será equiparada automaticamente à conversão para pastagem plantada. Também serão examinadas mudanças de classificação associadas à inundação, ao fogo e à detecção de usos mistos. Uma verificação de persistência do destino nos anos seguintes poderá reduzir transições espúrias, respeitando a perda de anos no final da série.

Até o momento, as bases processadas contêm estoques e participações de vegetação; as conversões brutas ainda precisam ser construídas. A diferença entre estoques anuais não será utilizada como substituto, pois combina supressão, regeneração, outros destinos e possíveis mudanças de classificação. Serão apresentadas, como sensibilidades, a área convertida em hectares e a conversão em proporção da área territorial fixa, de modo a avaliar o papel do denominador e da redução do estoque disponível.

\subsection{Preços, produção e capacidade produtiva}

A Produção Agrícola Municipal, tabela 1612, fornecerá área, produção e rendimento de soja e milho; a Pesquisa da Pecuária Municipal, tabela 3939, fornecerá o efetivo bovino \cite{IBGEPAM2026,IBGEPPM2026}. O rendimento observado não será confundido com capacidade produtiva exógena: ele também reflete clima, tecnologia, seleção das áreas cultivadas e decisões do produtor. Composição produtiva e produtividade históricas poderão definir exposições anteriores à amostra, após a revisão dos símbolos de ausência e zero do SIDRA. Áreas de soja e milho não serão somadas como se representassem terras físicas distintas, pois podem corresponder a safras sucessivas na mesma área.

As séries mensais da \textit{Pink Sheet}, do Banco Mundial, incluem preços internacionais de soja, milho e carne bovina \cite{WorldBank2026}. O arquivo local alcança agosto de 2026, mas o recorte empírico usará apenas observações compatíveis com a amostra. Soja e milho são cotados em dólares por tonelada, enquanto a carne bovina é cotada em dólares por quilograma. Variações logarítmicas serão calculadas por produto, após a escolha documentada da moeda, do deflator e da janela temporal. Preços internacionais não equivalem a preços recebidos pelo produtor em Mato Grosso; os dados locais de CEPEA e IMEA continuam pendentes. A inclusão da carne exigirá avaliar sua correspondência com a formação do preço local do boi.

A capacidade produtiva poderá ser aproximada por aptidão ou rendimento atingível do GAEZ v4, mantendo separadas as variantes de soja e milho, níveis de insumos e condições de sequeiro \cite{FAOGAEZ2021}. Os arquivos reunidos utilizam clima de referência de 1961--1990, mas a documentação do GAEZ v4 inclui uma base de terras de 2010. Logo, não serão considerados automaticamente predeterminados para 2001. Será preciso verificar quais camadas entram na medida escolhida e sua dependência de usos observados após o início da amostra. Na ausência dessa validação, o GAEZ ficará restrito à caracterização ou a uma análise temporalmente compatível. Uma exposição alternativa construída com informações anteriores a 2001 testará especialização histórica; não será apresentada como medida puramente biofísica de capacidade.

As planilhas de custos da CONAB poderão complementar a interpretação dos incentivos, desde que se consiga harmonizar localidades, tecnologias, unidades e conceitos de custo \cite{CONAB2026}. As células já extraídas não constituem, por si sós, uma série comparável de custos municipais. Sem essa harmonização, a combinação entre preços e capacidade será denominada exposição a incentivos de receita, preservando a distinção em relação a lucro ou retorno líquido esperado.

\subsection{Instituições, controles e descrição da amostra}

As medidas institucionais candidatas incluem autos de infração e informações sobre embargos do IBAMA \cite{IBAMA2026}. São registros de atuação administrativa federal; não constituem uma base de embargos judiciais nem abrangem integralmente a fiscalização estadual e municipal. Serão distinguidos eventos únicos, data de lavratura, tipo de infração, situação administrativa, localização e valor nominal registrado. Uma multa lavrada não equivale a uma multa paga. Da mesma forma, listas de embargos atualmente vigentes não permitem reconstruir automaticamente todas as restrições que existiram no passado.

A preparação local identificou divergências entre município informado e município obtido pelas coordenadas, além de insuficiência da série histórica de embargos. Essas alternativas de agregação serão reconciliadas e reportadas, sem transformar ausência de registros em zero antes de confirmar a cobertura. As contagens poderão ser expressas por área territorial ou estoque inicial fixo de vegetação, com denominadores documentados. O uso de uma escala fixa evita introduzir no indicador institucional a própria conversão contemporânea, mas não resolve sua endogeneidade.

As estatísticas descritivas mostrarão cobertura temporal, ausências, distribuição do estoque disponível, conversões por destino, exposições produtivas e registros de fiscalização. Serão apresentados mapas e séries por bioma e atividade, sem interpretar correlações como resultados causais. Características fixas de relevo e localização serão absorvidas pelos efeitos fixos territoriais nas regressões; controles que variam no tempo só serão incluídos quando sua posição no mecanismo econômico estiver clara. Clima local e mudanças de acesso a mercados serão avaliados como possíveis confundidores, enquanto crédito, área plantada e produção contemporâneos podem constituir respostas aos próprios incentivos e não integrarão automaticamente o conjunto de controles.

A construção do painel e a definição final da amostra dependerão, portanto, de três tarefas centrais: obter as conversões brutas, validar a geografia histórica e estabelecer medidas de exposição e instituições coerentes com a anterioridade temporal. As limitações descritas refletem a auditoria local das bases de setembro de 2026; não são resultados de modelos já estimados.

